\documentclass[border=10pt]{standalone}
% 'standalone' produces a tight bounding box around the drawing.
% 'border=10pt' adds a small margin around the content.

\usepackage{circuitikz}
% The main package for drawing electronic circuits in LaTeX.

\usepackage{xcolor}
% Lets us name colors for annotation text.

% ── Convenience color for annotation comments ──────────────
\definecolor{annotate}{RGB}{0,100,180}

\begin{document}


\section{Magnetic Buzzer}

\subsection{Overview}

The alert sounder is a passive magnetic buzzer driven by a pulse-width modulated (PWM)
signal from the Raspberry Pi Zero. A passive magnetic buzzer is preferred over an active
buzzer or piezoelectric transducer because its tone is warm and rounded rather than harsh,
and because both pitch and volume are controllable in software via the PWM frequency and
duty cycle respectively.

The buzzer operates from the \SI{5}{\volt} supply rail and draws approximately
\SI{30}{\milli\ampere} to \SI{40}{\milli\ampere} at full drive. Since this exceeds the
safe output current of a GPIO pin, and because the supply voltage exceeds the
\SI{3.3}{\volt} GPIO logic level, a small-signal NPN bipolar junction transistor is
used as a low-side switch between the buzzer and ground.

\subsection{Circuit Description}

The GPIO pin drives the base of the NPN transistor through a current-limiting resistor
R1. When the GPIO is high, the transistor saturates and connects the buzzer's negative
terminal to ground, completing the circuit through the \SI{5}{\volt} supply. When the
GPIO is low, the transistor is cut off and the buzzer is de-energised.

Because the buzzer is an inductive load, a flyback diode is placed across its terminals
to suppress the reverse voltage spike that occurs when the transistor switches off.
Without this diode the spike could damage the transistor.

Volume is controlled by adjusting the PWM duty cycle. A duty cycle of approximately
50\% produces full volume; reducing the duty cycle lowers the average power delivered
to the buzzer and therefore reduces loudness. Pitch is controlled by adjusting the PWM
frequency. The range of \SI{2000}{\hertz} to \SI{2500}{\hertz} produces a clear,
pleasant tone that is audible in both laboratory and outdoor environments. The buzzer
is rated at approximately \SI{85}{\decibel} to \SI{90}{\decibel} at \SI{10}{\centi\meter}.

The complete circuit is shown below.

\begin{verbatim}
    RPi GPIO (PWM) ---[ R1 1k ]---+--- MMBT3904 Base
                                  |
                             MMBT3904 Emitter --- GND

    +5V ---[ Buzzer CEM-1203(42) ]---+--- MMBT3904 Collector
                                     |
                               1N4148W (flyback)
                               Anode  -> Collector side
                               Cathode -> +5V side
\end{verbatim}

\subsection{Parts List}

\begin{table}[h]
\centering
\begin{tabular}{llll}
\hline
\textbf{Reference} & \textbf{Part Number} & \textbf{Description} & \textbf{Package} \\
\hline
BZ1 & CEM-1203(42)  & Passive magnetic buzzer, \SI{5}{\volt}, \SI{2048}{\hertz} rated & 12\,mm through-hole \\
D1  & 1N4148W       & Fast switching diode, flyback protection                        & SOD-123 \\
Q1  & MMBT3904      & NPN small-signal BJT, \SI{40}{\volt}, \SI{200}{\milli\ampere}  & SOT-23  \\
R1  & R1K-0805      & Resistor \SI{1}{\kilo\ohm}, base current limiter               & 0805    \\
\hline
\end{tabular}
\caption{Magnetic buzzer circuit components}
\label{tab:buzzer_parts}
\end{table}

\subsection{Programming Notes}

The buzzer is controlled using the RPi.GPIO software PWM facility. Any GPIO pin may
be used, but a hardware PWM pin (GPIO 18 or GPIO 12 on the Raspberry Pi Zero) is
preferred because software PWM can exhibit jitter under system load, which may cause
audible roughness in the tone.

Key relationships are:

\begin{itemize}
    \item \textbf{Frequency} sets pitch. The range \SI{1500}{\hertz} to
          \SI{3000}{\hertz} is recommended. Values below \SI{1000}{\hertz} may be too
          quiet for some magnetic buzzers, and values above \SI{4000}{\hertz} become
          increasingly harsh.
    \item \textbf{Duty cycle} sets volume. A value of 50 (50\%) gives maximum loudness.
          Values below approximately 10 may be too quiet to be useful. Values above 50
          do not increase loudness and may introduce distortion.
    \item \textbf{Duration} is implemented as a blocking \texttt{time.sleep()} call.
          Interruption logic is not implemented in this version.
\end{itemize}

The GPIO pin must be initialised as an output before the PWM object is created. The
PWM object should be started with a duty cycle of zero and frequency set before each
call to avoid a brief click at startup.

\subsection{Example Code}

\begin{verbatim}
import RPi.GPIO as GPIO
import time

BUZZER_PIN = 18  # Hardware PWM pin recommended

GPIO.setmode(GPIO.BCM)
GPIO.setup(BUZZER_PIN, GPIO.OUT)

_pwm = GPIO.PWM(BUZZER_PIN, 2000)
_pwm.start(0)


def sound_alert(frequency: int = 2000,
                volume: int = 50,
                duration: float = 1.0) -> None:
    """
    Sound the magnetic buzzer for a fixed duration.

    Parameters
    ----------
    frequency : int
        Tone frequency in Hz. Recommended range 1500-3000.
        Default 2000 Hz.
    volume : int
        Duty cycle as a percentage, 0-50. Values above 50
        do not increase loudness. Default 50 (full volume).
    duration : float
        Duration in seconds to sound the alert. Blocking.
        Default 1.0 second.
    """
    volume = max(0, min(50, volume))       # clamp to 0-50
    frequency = max(100, min(10000, frequency))  # clamp to safe range

    _pwm.ChangeFrequency(frequency)
    _pwm.ChangeDutyCycle(volume)
    time.sleep(duration)
    _pwm.ChangeDutyCycle(0)                # silence buzzer, leave PWM running
\end{verbatim}

\noindent To sound a two-second alert at \SI{2500}{\hertz} at three-quarter volume:

\begin{verbatim}
    sound_alert(frequency=2500, volume=37, duration=2.0)
\end{verbatim}

\noindent To clean up GPIO resources on program exit:

\begin{verbatim}
    _pwm.stop()
    GPIO.cleanup()
\end{verbatim}
